\documentclass[10pt,a4paper]{amsart}
\usepackage[margin=19mm]{geometry}
\usepackage{amsmath,amssymb,mathtools}
\usepackage[scr=rsfs,cal=euler]{mathalfa}
\usepackage{xfrac}
\usepackage[unicode,hidelinks,bookmarks=false]{hyperref}
\newcommand{\ie}{i.e.\ }
\newcommand{\cf}{cf.\ }
\newcommand{\resp}{respectively\ }
\newcommand{\Subsection}{\subsection}
\providecommand{\nobreakdash}{\nobreak\hbox{-}}
\setlength{\emergencystretch}{2em}
\multlinegap0pt
\usepackage{amssymb}
\usepackage[all]{xy}
\usepackage{tikz}
\newtheorem{lemma}{Lemma}
\newcommand{\bR}{\mathbb{R}}
\newcommand{\bfW}{\mathbf{W}}
\newcommand{\cL}{\mathcal{L}}
\newcommand{\del}{\partial}
\title{Sectorial covers over fanifolds}
\author{Hayato Morimura}
\date{Source: arXiv:2310.10084v3 (CC BY 4.0)}
\begin{document}
\maketitle

\noindent\emph{Source and licence.} Sectorial covers over fanifolds. Version arXiv:2310.10084v3 (CC BY 4.0), distributed by arXiv under the Creative Commons Attribution 4.0 International licence, http://creativecommons.org/licenses/by/4.0/, as recorded in the arXiv licence metadata for this exact version and shown on the abstract page https://arxiv.org/abs/2310.10084v3. Full licence notice: \texttt{licenses/arxiv-2310.10084-CC-BY-4.0.txt}.\par\smallskip

\noindent\textit{Editorial scope.} Counted unit: the article's lemma lem:W-sector (source0.tex lines 1721--1725). The article's complete original proof of it is delivered with it (source0.tex lines 1726--1895, one \texttt{proof} environment of 170 lines). No mathematics is written, repaired or re-derived: the statement and the proof are the article's own lines, in the article's own notation, and no retained line is substituted or re-flowed.  No further source-local context is needed: every cross-reference of every retained line resolves inside the two delivered spans.  Nothing was omitted from any retained span, so the package carries no omission record.  The retained lines cite 2 works, whose entries are transcribed verbatim from the article's own bibliography source in the e-print and delivered with short editorial labels recorded in the item's reference mapping.  No retained line contains a figure environment, an image-inclusion command or drawing code, so no image is delivered and the package declares no asset dependency.  Editorial formatting only, in addition: an AMS \texttt{amsart} preamble that restates the article's own declarations and macros used by the retained lines, namely the environments \texttt{lemma} on separate counters, together with the standard packages those lines need.  The retained environments print in this document as Lemma 1 in order of appearance, whereas the article prints the same items with the numbering of its own full document.  The complete native source of the article as distributed in the arXiv e-print for this exact version is bundled unchanged as \texttt{sources/148514/source0.tex}.
\begin{lemma} \label{lem:W-sector}
Each
$\widetilde{\bfW}(P)$
is a Weinstein sector.
\end{lemma}

\begin{proof}
We adapt the idea in
\cite[Example 1.33]{GPS2}.
Recall that
we are assuming
$\Phi$
to be closed.
For simplicity,
we will further assume that
there is no lower dimensional strata
which are not adjacent to higher dimensional strata.
Since by construction
$\bar{\pi}(\del^2 \widetilde{\bfW}(P))$
is away from
$P$,
we may assume that
$\del^2 \widetilde{\bfW}(P)$
is disjoint from
$T^* \widehat{M}_P
\subset
\widetilde{\bfW}(P)
\cap
\widetilde{\bfW}(\Phi_0)$.
Let
$I$
be a $1$-stratum
which intersects
$\bar{\pi}(\del^2 \widetilde{\bfW}(P))$
at a point
$p$.
Due to the above perturbation,
locally in
$\Phi_1$
the image
$\bar{\pi}(\del^2 \widetilde{\bfW}(P))$
becomes perpendicular to
$I$.
Since
$\bar{\pi}$
kills
the base direction of
$T^* \widehat{M}_I$
and
the fiber direction of
$T^* I_\circ$,
over the neighborhood of
$I$
the boundary
$\del^2 \widetilde{\bfW}(P)$
is isomorphic to a product of
$T^* \widehat{M}_I$
and
the fiber
$T^*_p I_\circ$.
Inductively,
let
$S^{(k)}$
be a $k$-stratum adjacent to a $(k-1)$-stratum
$S^{(k-1)}$
which intersects
$\bar{\pi}(\del^2 \widetilde{\bfW}(P))$.
Due to the above perturbation,
locally in
$\Phi_k$
the image
$\bar{\pi}(\del^2 \widetilde{\bfW}(P))$
becomes perpendicular to
$S^{(k)}$.
Since
$\bar{\pi}$
kills
the base direction of
$T^* \widehat{M}_{S^{(k)}}$
and
the fiber direction of
$T^* S^{(k)}_\circ$,
over the neighborhood of
$S^{(k)}$
the boundary
$\del^2 \widetilde{\bfW}(P)$
is isomorphic to a product of
$T^* \widehat{M}_{S^{(k)}}$
and
the restriction
$T^* S^{(k)}_\circ |_{\bar{\pi}(\del^2 \widetilde{\bfW}(P)) \cap S^{(k)}}$.
By construction locally in
$\Phi_{(k-1)}$
the intersection
$\bar{\pi}(\del^2 \widetilde{\bfW}(P)) \cap S^{(k)}$
becomes perpendicular to
$S^{(k-1)}$,
the handle attachment
$\bfW(\Phi_{k-1}) \#_{\cL_{k-1}} (T^* \widehat{M}_{S^{(k)}} \times T^* S^{(k)}_\circ)$
glues this product to the gluing of the products in the previous steps.

Rewinding the process,
one sees that
$\del^2 \widetilde{\bfW}(P)$
is obtained by extending via the handle attachments the union of the inverse images under
$\bar{\pi}$
of hypersurfaces in the intersection of
$\Phi_n \setminus \Phi_{n-1}$
and
the connected component containing
$P$.
We claim that
this hypersurface
$\del^2 \widetilde{\bfW}(P)$
is sectorial.
Since the newly attached pieces are cylindrical,
so is
$\del^2 \widetilde{\bfW}(P)$.
It remains to show the existence of an $\alpha$-defining function on
$\del^2 \widetilde{\bfW}(P)$.  
By
\cite[Construction 12.17]{GPS2}
we may assume
$\del^2 \widetilde{\bfW}(P)$
to be one of the hypersurfaces-with-corners dividing
$\widetilde{\bfW}(\Phi)$
into the original connected components.
Let
$I$
be a $1$-stratum adjacent to
$P$
and
$H_{P, I}$
any of such hypersurfaces intersecting
$I$.
Fix the standard coordinates
$(x, y)$
on the symplectic manifold
$T^* I_\circ$.
Unwinding
the filtration of
$\Phi$
and
the construction of
$\widetilde{\bfW}(\Phi)$,
one sees that
$H_{P, I}$
is given by
$\{ x = 0 \}$
up to perturbation.
We denote by
$y_{P, I}$
the function on a neighborhood of
$H_{P, I}$
which is canonically induced by
$y \colon T^* I_\circ \to \bR$.
Clearly,
$y_{P, I}$
is linear near infinity.
The Hamiltonian vector field
$X_{y_{P, I}}$
is outward pointing along
$H_{P, I}$,
as it is given by a lift of
$\del_x$
on 
$T^* I_\circ$.
Hence
$H_{P, I}$
is sectorial.
Thus the Weinstein manifold-with-boundary
$\widetilde{\bfW}(P)$
satisfies the condition in
\cite[Definition 2.4]{GPS1}
to be a Liouville sector.
\end{proof}

\begin{thebibliography}{99}
\bibitem[GPS1]{GPS1} S. Ganatra, J. Pardon, and V. Shende, \emph{Covariantly functorial wrapped Floer theory on Liouville sectors}, Publ. math. IHES. 131, 73–200 (2020).
\bibitem[GPS2]{GPS2} S. Ganatra, J. Pardon, and V. Shende, \emph{Sectorial descent for wrapped Fukaya categories}, J. Amer. Math. Soc. 37, 499-635 (2024).
\end{thebibliography}
\end{document}
